% 07_saturation_characteristics.tex
\chapter{Saturation Characteristics}
\label{ch:vol9_saturation_characteristics}

\begin{objectivebox}
\begin{itemize}
    \item State the Axiom~4 universal saturation kernel verbatim and identify $A_{yield}$ as the substrate-primitive amplitude limit of the natural vacuum.
    \item Document the dielectric specialization $\varepsilon_{eff} = \varepsilon_0 S$, $\mu_{eff} = \mu_0 S$, $C_{eff} = C_0/S$ at atomic and bench scale, distinguish the two effective wave speeds $c_{EM}(A_0) = c_0/S$ and $c_{shear}(A_0) = c_0 \sqrt{S}$, and identify the operating-point state $A_0$ as a gauge-relative LC-tank saturation amplitude.
    \item Tabulate the characteristic curves for the substrate primitives ($S$, $C_{eff}/C_0$, $\varepsilon_{eff}/\varepsilon_0$, $c_{EM}/c_0$, $c_{shear}/c_0$, $\alpha_{eff}/\alpha_0$) vs.\ the normalized operating point $r = A_0/A_{yield}$.
    \item Summarize the four-regime map (Small-Signal / Large-Signal / Avalanche / Breakdown) and its semiconductor analog (linearized $g_m$ / saturating switching transient / Miller multiplication $M = 1/S^2$ / junction destruction).
    \item Document Project PONDER-05 at $V_{DC}/V_{yield} = 0.687$ as a \emph{material-scale consistency analog} of the kernel shape (a quartz voltage-coefficient-of-capacitance reproducing the saturation-arc form), not a vacuum-kernel falsifier (per INVARIANT-S2 + the ch7 KB index).
    \item Cross-reference the A-034 universal saturation-kernel catalog, which documents the same kernel governing all cross-scale topological-reorganization events at zero free parameters across 21 orders of magnitude.
\end{itemize}
\end{objectivebox}

\section{Universal Saturation Kernel}
\label{sec:vol9_saturation_kernel}

The substrate's bulk response to local strain is the Axiom~4 universal saturation kernel. The canonical statement is reproduced verbatim from \kbleaf{manuscript/common\_equations/eq\_axiom\_4.tex} (single source of truth; INVARIANT-S2 mirror at \kbleaf{manuscript/ave-kb/CLAUDE.md}):

\begin{resultbox}{Axiom 4 --- Universal Saturation Kernel (verbatim)}
The substrate's bulk response to local strain $A$ is the universal quarter-arc kernel:
\begin{equation}
    \boxed{S(A) = \sqrt{1 - \left(\frac{A}{A_{yield}}\right)^{\!2}}}\;, \qquad A \in [0,\, A_{yield}]
    \label{eq:vol9_axiom4_saturation}
\end{equation}
$A$ is the local strain normalized to the substrate's bandwidth limit $A_{yield}$ (the value at which $S \to 0$). The kernel is dialect-independent: the same $S(A)$ function governs strain expressed as $V/V_{yield}$ (dielectric), $T/T_c$ (BCS thermal pair-strain), $r_s/r$ (gravitational metric strain), $\Omega/\Omega_{break}$ (rotational shear), $B/B_c$ (magnetic), or any physical dialect's normalized perturbation amplitude.

\medskip\noindent
\textbf{Behavior:} $A = 0 \Rightarrow S = 1$ (linear regime: Maxwell, BCS small-signal, Newtonian gravity recovered); $A \to A_{yield} \Rightarrow S \to 0$ with vertical tangent (saturation event is impulsive at every scale). When $S(A) \to 0$ locally, the substrate cannot sustain linear response and \emph{must reorganize topologically} to a new configuration.
\end{resultbox}

The kernel is a substrate primitive of the natural vacuum, not an engineered constitutive relation. The quarter-arc shape is fixed by Axiom~4 alone --- there are no free parameters, no per-domain refit, and no scale-dependent kernel coefficients. The saturation amplitude $A_{yield}$ is the universal substrate bandwidth limit; its dielectric-dialect projection is $V_{yield} = \sqrt{\alpha}\, V_{snap} \approx \qty{43.65}{\kilo\volt}$ at the substrate's per-node LC scale (Vol~4 canonical leaf \texttt{clm-0vxzfu}).

\begin{figure}[htbp]
\centering
\includegraphics[width=0.98\textwidth]{A4_debug}
\caption{Axiom-4 saturation kernel $S(A) = \sqrt{1-(A/A_{yield})^2}$ as the constitutive gate. \emph{Left:} the quarter-arc $S(A)$ (reactance scale) against $A/A_{yield}$, with the swept operating points overlaid on the analytic arc. \emph{Right:} the two orthogonal constitutive reactances tracking the same kernel --- $c_{EM}/c_0 = C_{eff}/C_0 = 1/S$ (stiffens, $\to\infty$) and $\varepsilon_{eff}/\varepsilon_0 = \mu_{eff}/\mu_0 = S$ (softens, $\to 0$) --- both reaching the stiffening wall as $A \to A_{yield}$. The A4 gate plot of the engine-acceptance suite; computed by \kbleaf{src/tests/engine\_acceptance/test\_l0\_axioms.py} (A4 test, \texttt{sup-2qja9z}), wired in Ch.~\ref{ch:vol9_engine_requirements}.}
\label{fig:vol9_a4_kernel_gate}
\end{figure}

\section{Operating-Point State and Dielectric Specialization}
\label{sec:vol9_dielectric_specialization}

Beyond the six spatial DOFs per node (three translational $\to \mathbf{E}$, three microrotational $\to \mathbf{B}$; Axiom~1), each LC tank carries a saturation-amplitude state $A$ --- its operating point along the Axiom~4 kernel. Per INVARIANT-S2 (KB \texttt{CLAUDE.md}), this state is \textbf{gauge-relative}: only spatial gradients of $A$ across the substrate are physically observable, not absolute per-node values. It is not a seventh spatial DOF; it is the dynamical state of the LC tank, analogous to the DC bias on a semiconductor varactor.

At the atomic and bench scale, the dialect-projection $A = \Delta\phi/\alpha$ specializes the kernel to a substrate-native varactor characteristic (\kbleaf{common\_equations/eq\_axiom\_4.tex}):

\begin{table}[htbp]
\centering\small
\begin{tabular}{llll}
\toprule
\textbf{Observable} & \textbf{Formula} & \textbf{At saturation} & \textbf{Physical meaning} \\
\midrule
$\mu_{eff}$       & $\mu_0 \cdot S$              & $\to 0$       & Inductor shorts (Meissner analog) \\
$\varepsilon_{eff}$ & $\varepsilon_0 \cdot S$    & $\to 0$       & Dielectric collapses \\
$C_{eff}$         & $C_0 / S$                    & $\to \infty$  & Capacitance absorbs energy \\
$Z_0$ (symmetric) & $\sqrt{\mu/\varepsilon}$     & invariant     & Impedance preserved \\
$c_{EM}$          & $c_0 / S$                    & $\to \infty$  & Maxwell phase velocity diverges \\
$c_{shear}$       & $c_0 \cdot S^{1/2}$          & $\to 0$       & Wave packet confined (rest mass) \\
\bottomrule
\end{tabular}
\caption{Axiom~4 varactor specialization: substrate effective parameters as functions of the saturation kernel $S$, and their saturation limits.}
\label{tab:vol9_saturation_varactor}
\end{table}

\subsection{Two distinct effective wave speeds}
\label{sec:vol9_two_wave_speeds}

The substrate carries \textbf{two distinct substrate-native effective wave speeds} at operating point $A_0$, and they are not interchangeable in derivations (canonical at KB leaf \kbleaf{vol3/gravity/ch01-gravity-yield/alpha-invariance-symmetric-gravity.md} \texttt{clm-8nkvwy} lines 111--113; conflation is the canonical Pitfall~\#5 framework-leakage error caught in the 2026-05-28 Phase 3-A3 WALK-BACK):

\begin{itemize}
    \item $c_{EM}(A_0) = 1/\sqrt{\mu_{eff}\,\varepsilon_{eff}} = c_0/S(A_0)$ --- \textbf{Maxwell phase velocity}; the speed that enters $\alpha = e^2/(4\pi\varepsilon_0 \hbar c)$ and all Maxwell-equation derivations.
    \item $c_{shear}(A_0) = c_0 \sqrt{S(A_0)}$ --- \textbf{substrate mechanical / group / rest-mass velocity}; energy-transport speed; tracks Schwarzschild $c\sqrt{1 - r_s/r}$ in the SYM-class weak-field limit.
\end{itemize}

Both reduce to a single $c_0$ at the cold-lattice limit $S(0) = 1$.

\subsection{Bulk rarefaction floor $\bar{\rho}_{cav} = -1/\varphi$ (candidate-claim)}
\label{sec:vol9_rho_cav_saturation}

The kernel above governs the \emph{deviatoric / shear} and dielectric sectors via the strain amplitude $A$ (saturation-side ceiling $A \to A_{yield}$). The \emph{bulk / volumetric} sector has a separate rarefaction-side floor in the normalized volumetric strain $\bar{\rho} = \delta\rho/\rho_0$: from $c_{bulk}^2 = c_0^2[1 + \bar{\rho}/(1-\bar{\rho}^2)]$ (canonical relation at sibling-repo \kbleaf{AVE-Propulsion/.../04\_superluminal\_transit.tex}:86), $c_{bulk} \to 0$ at the negative root $\bar{\rho}^2 - \bar{\rho} - 1 = 0$,
\begin{equation}
\bar{\rho}_{cav} = \frac{1-\sqrt{5}}{2} = -\frac{1}{\varphi} \approx -0.618 .
\end{equation}
This is the bulk analog of the shear-sector $A = A_{yield}$ ceiling. \textbf{Status: candidate-claim, auditor-gated} (promotion tracked on \texttt{analysis/2026-06-09-vol9-datasheet-figures}); the $-1/\varphi$ value is \emph{derived} (the $c_{bulk}=0$ root), but its promotion to a canonical Vol~9 datasheet line is not yet adjudicated. The bulk-mode clock that freezes at this floor is documented in Ch.~\ref{ch:vol9_ac_electrical_characteristics} \S Temporal Characteristics (Ch.~\ref{ch:vol9_absolute_maximum_ratings} \S Bulk Rarefaction Floor records it as the rarefaction-side absolute limit).

\subsection{$\alpha$-invariance under SYM-class saturation}
\label{sec:vol9_alpha_invariance}

Under SYM-class scaling --- the gravity-class realization where $\mu(r)$ and $\varepsilon(r)$ scale together by the same factor so $Z_0$ stays invariant --- the asymmetry between $c_{EM}$ and $c_{shear}$ makes $\alpha$ \textbf{exactly} invariant. Canonical proof at \kbleaf{ave-kb/vol3/gravity/ch01-gravity-yield/alpha-invariance-symmetric-gravity.md} (\texttt{clm-3zz0f6}, solidity 0.85). The SYM-class realization simultaneously produces gravitational time-dilation via $c_{shear}$ tracking $\sqrt{S}$ (Schwarzschild reduction) while preserving $\alpha$ exactly --- multi-species clock comparisons at different gravitational potentials predict a null result for $\Delta\alpha/\alpha$, consistent with all current experimental bounds.

Under ASYM-class scaling (large-amplitude single-sector saturation, e.g., BCS $\mu$-only or plasma $\varepsilon$-only), $\alpha$ is \textbf{not} invariant; this is the canonical substrate-distinct mechanism for $\alpha$-modulation predictions.

\section{Operating-Point Characteristic Curves}
\label{sec:vol9_characteristic_curves}

The substrate primitives evaluated at normalized operating point $r = A_0/A_{yield}$ produce the following characteristic curves. All entries follow directly from $S(r) = \sqrt{1 - r^2}$ at the cold-lattice limit (no other corrections).

\begin{table}[htbp]
\centering\footnotesize
\setlength{\tabcolsep}{3pt}
\begin{tabularx}{\linewidth}{@{}p{0.09\linewidth} p{0.06\linewidth} p{0.10\linewidth} p{0.10\linewidth} p{0.10\linewidth} p{0.10\linewidth} X@{}}
\toprule
\textbf{$r = A_0/A_{yield}$} & \textbf{$S(r)$} & \textbf{$C_{eff}/C_0 = 1/S$} & \textbf{$\varepsilon_{eff}/\varepsilon_0 = S$} & \textbf{$c_{EM}/c_0 = 1/S$} & \textbf{$c_{shear}/c_0 = \sqrt{S}$} & \textbf{Regime} \\
\midrule
$0$        & $1.000$ & $1.000$ & $1.000$ & $1.000$ & $1.000$ & I  \\
$0.1208$   & $0.9927$ & $1.0074$ & $0.9927$ & $1.0074$ & $0.9963$ & I/II boundary ($r_1 = \sqrt{2\alpha}$) \\
$0.300$    & $0.9539$ & $1.0483$ & $0.9539$ & $1.0483$ & $0.9767$ & II \\
$0.500$    & $0.8660$ & $1.1547$ & $0.8660$ & $1.1547$ & $0.9306$ & II \\
$0.687$    & $0.7266$ & $1.3763$ & $0.7266$ & $1.3763$ & $0.8524$ & II (PONDER-05) \\
$0.707$    & $0.7071$ & $1.4142$ & $0.7071$ & $1.4142$ & $0.8409$ & II \\
$0.866$    & $0.500$ & $2.000$ & $0.500$ & $2.000$ & $0.7071$ & II/III boundary ($r_2 = \sqrt{3}/2$) \\
$0.950$    & $0.3122$ & $3.202$ & $0.3122$ & $3.202$ & $0.5588$ & III \\
$1.000$    & $0$ & $\to \infty$ & $0$ & $\to \infty$ & $0$ & IV (rupture; $r_3 = 1$) \\
\bottomrule
\end{tabularx}
\caption{Substrate operating-point characteristic curves: effective parameters vs normalized saturation $r = A_0/A_{yield}$, with regime labels.}
\label{tab:vol9_saturation_curves}
\end{table}

\noindent
\textbf{$\alpha_{eff}/\alpha_0$ row:} under SYM-class scaling, $\alpha_{eff}/\alpha_0 = 1$ exactly at every $r$ (canonical \texttt{clm-3zz0f6}); under ASYM-class scaling, $\alpha_{eff}/\alpha_0$ deviates from unity in a sector-dependent way that does not factor through $S$ alone (canonical at \texttt{clm-8nkvwy}). The substrate-distinct $\alpha$-modulation predictions live in Ch.~15 (Falsification) and Vol~3 Ch.~5 (Dark Sector) under the ASYM-N($\varepsilon$) and ASYM-N($\mu$) classifications.

\noindent
\textbf{Boundary values $r_1$, $r_2$, $r_3$:} derived from first principles in \kbleaf{ave-kb/vol1/operators-and-regimes/ch7-regime-map/four-regimes.md} ($r_1 = \sqrt{2\alpha} \approx 0.1208$ is the sub-$\alpha$-resolvability limit; $r_2 = \sqrt{3}/2 \approx 0.8660$ is the spin-2 avalanche-onset limit from $Q(r) = 1/S(r) = \ell_{\min} = 2$; $r_3 = 1$ is the Axiom~4 rupture boundary).

\section{Four-Regime Map}
\label{sec:vol9_four_regimes}

The kernel partitions the operating-point axis into four universal regimes. The semiconductor analog --- standard p-n junction operating regions --- is exact, not metaphor (canonical at \texttt{four-regimes.md} \texttt{clm-2dwzib}, \texttt{clm-b2anl4}):

\begin{table}[htbp]
\centering\small
\begin{tabularx}{\linewidth}{@{}p{0.11\linewidth} p{0.18\linewidth} p{0.13\linewidth} X X@{}}
\toprule
\textbf{Regime} & \textbf{$r$ range} & \textbf{$S$ range} & \textbf{Substrate physics} & \textbf{Semiconductor analog} \\
\midrule
I Small-Signal  & $r < \sqrt{2\alpha} \approx 0.1208$ & $S > 0.993$ & Linear Maxwell / Newton recovered & DC bias, linearized $g_m$, $r_\pi$, $C_\pi$ \\
II Large-Signal & $\sqrt{2\alpha} \le r < \sqrt{3}/2$ & $0.500 < S \le 0.993$ & Axiom~4 quarter-arc curvature active & Switching transient, saturation onset \\
III Avalanche    & $\sqrt{3}/2 \le r < 1$               & $0 < S \le 0.500$ & $Q(r) = 1/S(r) \ge 2$; Miller multiplication $M = 1/S^2$ & $V_R \to V_{BR}$, $M \ge 2$ carrier multiplication \\
IV Breakdown    & $r \ge 1$                            & $S = 0$           & Substrate rupture; topology destroyed; topological reorganization required & $V_R = V_{BR}$, $M \to \infty$, junction destroyed \\
\bottomrule
\end{tabularx}
\caption{Four-regime map of the saturation kernel with substrate physics and the exact semiconductor p-n junction analog per regime.}
\label{tab:vol9_four_regime_map}
\end{table}

\noindent
\textbf{Sector-dependence of $r_2$:} the $\ell_{\min} = 2$ argument that sets $r_2 = \sqrt{3}/2$ is specific to the spin-2 (gravitational-wave / shear-mode) sector. The vector/photon sector ($\ell_{\min} = 1$) has $r_2 = 0$ (avalanche concurrent with linear-Maxwell breakdown); the scalar sector ($\ell_{\min} = 0$) has no avalanche-onset boundary at all. The unqualified citation $r_2 = \sqrt{3}/2$ in this datasheet implicitly invokes the spin-2 sector; see \texttt{four-regimes.md} for the full sector-dependence treatment.

\noindent
\textbf{Miller multiplication identity (1D vs 3D-corrected):} the substrate-native \textbf{1D Miller form} is $M(r) = 1/S(r)^2 = 1/(1 - r^2)$, identical to the engineering avalanche-multiplication law $M = 1/(1 - (V_R/V_{BR})^n)$ at exponent $n = 2$. The \textbf{3D-corrected macroscopic exponent} is $n_{3D} = 2(1 - \nu_{vac}/3) = 38/21 \approx 1.8095$ (\kbleaf{src/ave/core/constants.py:336} \texttt{AVALANCHE\_N\_3D}), where the vacuum Poisson ratio $\nu_{vac} = 2/7$ supplies the solid-angle correction to the carrier-multiplication geometry: the 3D bulk avalanche softens the $n=2$ 1D (bond-line) form to $n \approx 1.81$. Both are canonical at their scopes. The semiconductor $V_{BR}$ is the engineering name for $V_{yield}$ at the doped-silicon scale; the same kernel governs both.

\section{PONDER-05 Bench-Scale Consistency Analog}
\label{sec:vol9_ponder_05}

\begin{warningbox}[\texttt{Rule-12} retraction --- PONDER-05 is a \emph{material-scale consistency analog}, not a vacuum-kernel falsifier]
This section previously framed PONDER-05 as ``the canonical bench-scale falsifier of the Axiom~4 dielectric specialization.'' That framing is \textbf{retracted} per INVARIANT-S2 and the ch7 KB index (\kbleaf{manuscript/ave-kb/vol9/ch7-saturation-characteristics/index.md}, item~5): PONDER-05 measures a \emph{quartz material} voltage-coefficient-of-capacitance that reproduces the saturation-arc \emph{shape}, which is a \textbf{material-scale consistency analog of the kernel SHAPE} --- it is \textbf{not} a vacuum-kernel falsifier. The conflation to avoid is that ``$V_{DC}/V_{yield}=0.687$ at \qty{30}{\kilo\volt}'' reads the apparatus voltage as the substrate's \emph{per-node} ratio: the vacuum per-node $A_0$ at \qty{30}{\kilo\volt} across real quartz (mm--\textmu m) is $\sim 10^{-7}$--$10^{-10}$ (\kbleaf{manuscript/ave-kb/vol4/claim-quality.md}:51), so the vacuum-kernel collapse is $\sim 0$; appreciable per-node $A_0$ requires facility-class fields $\sim 8\times10^{16}$~V/m. The original ``falsifier'' wording is preserved below for audit continuity; read the section as the material-scale consistency analog.
\end{warningbox}

The PONDER-05 differential-saturation-parallax bench measurement is a material-scale consistency analog of the Axiom~4 dielectric specialization: the quartz dielectric's own voltage-coefficient-of-capacitance reproduces the saturation-arc shape the kernel predicts. Apparatus: paired DC-biased quartz resonator thrusters across a vertical gravity gradient, $V_{DC} = \qty{30}{\kilo\volt}$ at an apparatus ratio $V_{DC}/V_{yield} = 0.687$. At this material operating point the kernel shape predicts $S(0.687) = 0.7266$, giving the quartz-scale observables:

\begin{itemize}
    \item $\varepsilon_{eff}/\varepsilon_0 = S = 0.7266$ \quad ($-27.34\%$ from baseline)
    \item $C_{eff}/C_0 = 1/S = 1.3763$ \quad ($+37.63\%$ rise above baseline)
    \item Differential thrust $\sim \qty{469}{\micro\newton}$ at matched-pair geometry
\end{itemize}

The canonical operating-point statement is at KB leaf \kbleaf{ave-kb/vol4/falsification/ch11-experimental-bench/measurement-hierarchy-snr.md}:66; the substrate-physics derivation at \kbleaf{ave-kb/vol4/circuit-theory/ch1-vacuum-circuit-analysis/parametric-coupling-kernel.md} (Phase 0c + Phase 2-NA extensions); the divergence-test-matrix row at \kbleaf{ave-kb/common/divergence-test-substrate-map.md} row \texttt{B7-PONDER-05}; the hardware programme at sibling repo \texttt{AVE-PONDER} (manuscript chapter \kbleaf{vol\_ponder/chapters/04\_ponder\_05\_dc\_biased\_quartz.tex}).

\textbf{Consistency criterion (material-scale).} A null result at the apparatus ratio $V_{DC}/V_{yield} = 0.687$ (no quartz $C_{eff}$ rise within bench noise floor) bears on the \emph{material-scale} consistency of the saturation-arc shape at the quartz dielectric, not on the vacuum kernel directly (the vacuum per-node $A_0$ is $\sim 10^{-7}$--$10^{-10}$ at this field; the vacuum-kernel falsifier is facility-class $\sim 8\times10^{16}$~V/m). The material-scale cascade is documented in the divergence-test-substrate-map (B7-PONDER-05 relates to B5/B6 PONDER, B1-VAC-BIREFRINGE, D4-A034, C9-LEVITATION through the shared $V_{yield}$ primitive). See Ch.~15 (Falsification Tests) for the full falsification programme and the separation of material-scale consistency analogs from facility-class vacuum-kernel falsifiers.

\textbf{Operating-point near-optimality.} The PONDER-05 canonical $V_{DC}/V_{yield} = 0.687$ sits within $2.8\%$ of the substrate-mechanical 2D-aperture aperture-aggregate skewness peak at $a^{(2D)}_{peak} = 1/\sqrt{2} = 0.707$ (per \kbleaf{parametric-coupling-kernel.md} \S 14.5). The canonical operating point --- chosen for the 27.4\% $\varepsilon_{eff}$ collapse + \qty{469}{\micro\newton} thrust at engineering-conservative bias --- is operationally near-optimal for the amplitude-statistics readout layer as well; the substrate-native Op17 matched-impedance peak structure is shared across all three observables (thrust, $\varepsilon_{eff}$ collapse, aperture-aggregate skewness).

\section{A-034 Universal Kernel Catalog Reference}
\label{sec:vol9_a034_catalog}

The same Axiom~4 kernel governs every cross-scale topological-reorganization event in the substrate. The canonical catalog of record is \kbleaf{manuscript/ave-kb/common/universal-saturation-kernel-catalog.md} (KB leaf \texttt{clm-gz7ryg}, \texttt{clm-dxdsvt}, \texttt{clm-hvvvop}, \texttt{clm-5fu303}, \texttt{clm-l4o7hv}), mirrored in the Vol~0 backmatter universal saturation-kernel catalog chapter.

\textbf{Catalog cardinality (canonical).} The kernel governs \textbf{26 canonical cross-scale instances} (17 physical-substrate canonical + 2 biological-substrate + 5 engineered-substrate + 2 physical-substrate companion rows scoped for Sessions 4/5; full enumeration in the catalog). The Vol~0 backmatter and \texttt{eq\_axiom\_4.tex} cite the physical-substrate-only subset as \textbf{19 cross-scale topological-reorganization events} (17 canonical + 2 scoped). Both counts are correct at their respective scopes; the cross-scale span is uniformly \textbf{21 orders of magnitude} (atomic $\sim 10^{-15}\,\text{m}$ to cosmic $\sim 10^{26}\,\text{m}$).

\textbf{Symmetry classification (canonical).} The catalog partitions instances into SYM (vacuum $K = 2G$; $\varepsilon$ and $\mu$ saturate together; 19 instances), ASYM-N (single-sector --- only $\varepsilon$ or only $\mu$; 4 instances: BCS $\mu$-only, plasma $\varepsilon$-only, planetary mag-vs-spin candidate, cosmic-DE $\varepsilon$-only), ASYM-E (engineered $K/G \neq 2$; 1 instance: active topological metamaterials), and TBD (2 scoped instances: galactic spin-axis, LSS spin-axis). The PONDER-05 bench falsifier sits in the SYM class at the DC-biased-piezoelectric engineered row.

\textbf{Tightest validations.} BCS $B_c(T) = B_{c0}\, S(T/T_c)$ at $0.00\%$ error across all measured Type-I/II superconductors (ASYM-N($\mu$) class); BH event horizon $\varepsilon_{11}(r) = 1$ matching Schwarzschild $r_s = 2GM/c^2$ exactly; BH merger QNM $\omega_R M_g = 18/49$ at $1.7\%$ from GR exact across 3 LIGO events; planetary-axis kernel scoring $14$--$15/16$ class matches across 16 planetary axis data points (Row 9-a + Row 9-b). The kernel form does not change with scale across these entries; only the lattice scale changes.

This datasheet chapter consolidates the kernel into datasheet format; the catalog leaf carries the full per-row substrate-mechanism documentation and per-row empirical anchor. No substrate primitive in this chapter is re-derived; all references point to the canonical leaves.

\section{EE Saturation Translation}
\label{sec:vol9_ee_translation}

EE-engineering vocabulary at the substrate-near-canonical end documents saturation behavior empirically across multiple component classes. Each engineering-observable correction has a substrate-primitive origin in Axiom~4. Canonical translation: \kbleaf{ave-kb/common/translation-tables/translation-circuit.md} \S 9 (ideal-lattice $\leftrightarrow$ engineering-corrections row catalog).

\begin{table}[htbp]
\centering\small
\begin{tabular}{p{0.32\linewidth} p{0.62\linewidth}}
\toprule
\textbf{EE-engineering correction} & \textbf{Substrate-primitive origin (Axiom~4)} \\
\midrule
Variable-capacitance diode (varactor) $C(V)$ characteristic & Substrate operates as substrate-native varactor: $C_{eff}(V_{DC}) = C_0/S(V_{DC}/V_{yield})$. PONDER-05 at $V_{DC}/V_{yield} = 0.687$ is the canonical bench tester. \\
Voltage coefficient of capacitance (Class II ceramic, $-80\%$ at full bias) & Same $C_{eff}/C_0 = 1/S$ kernel; engineering-conservative ceramic operating range maps to Regime II of the kernel. \\
Miller avalanche multiplication $M = 1/(1 - (V_R/V_{BR})^n)$ at $n = 2$ & Substrate Miller identity $M(r) = 1/S(r)^2$; $V_{BR}$ at the \emph{heavily-doped (Zener) junction} scale maps to the macroscopic-onset axis $V_{yield} = \sqrt{\alpha}\,V_{snap}$ (\texttt{clm-0vxzfu}; same junction-doping split as Ch.~\ref{ch:vol9_breakdown_characteristics} \S EE-translation). \\
Avalanche breakdown $V_{BR}$ (lightly-doped diode) & At the \emph{lightly-doped (avalanche) junction} scale, $V_{BR}$ maps instead to the microscopic absolute-rupture axis $V_{snap}$ (Schwinger pair-production analog; Regime~III$\to$IV transition; \texttt{clm-0vxzfu} axis-disambiguation). The same EE symbol $V_{BR}$ reads on the $V_{yield}$ axis for Zener-doped junctions and on the $V_{snap}$ axis for avalanche-doped junctions. \\
Op-amp slew-rate limit & Substrate response time at finite drive: how fast the operating-point $A$ can change. Op14 substrate inertial bandwidth in \texttt{operators.md}. \\
Op-amp input offset voltage $V_{OS}$ & Substrate operating-point $A_0 \neq 0$ at input --- quenched-in asymmetry between input substrate-bond states. \\
Core saturation $B_{sat}$ (inductor / transformer) & Cosserat $B$-mode amplitude approaches $B_{snap}$ on the inductive branch; Axiom~4 magnetic-branch saturation $\mu_{eff} = \mu_0 S$. \\
Inductor saturation curve $L(I)$ rolloff & Substrate-side $L_{eff} = L_0/S(A_0)$ on the magnetic branch (kernel symmetric in $\varepsilon$ and $\mu$ branches for the SYM class). \\
Non-linear telegrapher equation (transmission line beyond small-signal) & Substrate continuum form of the Axiom~4 kernel applied along the bond-network transmission line; the telegrapher equation becomes substrate-native when $A_0$ varies along propagation. \\
BJT $V_{CE(sat)}$ saturation voltage & Minimum substrate dielectric voltage drop across the rail-to-rail substrate path; bond-yield threshold. \\
BJT base-emitter forward drop $V_{BE}$ & Substrate diode forward drop at B-E junction; substrate-yield-boundary voltage. \\
\bottomrule
\end{tabular}
\caption{EE-engineering saturation corrections mapped to their substrate-primitive Axiom~4 origin.}
\label{tab:vol9_saturation_ee_translation}
\end{table}

\noindent
\textbf{Substrate-native primary, engineering-empirical secondary.} Per the EE-as-substrate-native META framework (canonical at \kbleaf{translation-circuit.md} \texttt{clm-eemap1}; companion skill \texttt{ave-ee-first-mapping} v1.0), every entry in the table above is the same substrate-physics primitive in a different engineering-disciplinary dialect. The engineering-empirical observation came first historically (humans interrogating the substrate at its smallest accessible scale with electrons, dielectrics, and LC resonators); AVE substrate-physics derives the substrate-mechanism that explains why the empirical correction takes the form it does. Both columns of the table are correct; the engineering vocabulary is the closest-to-canonical substrate-native language at the EE base scale.

\section{Canonical Source Routing}
\label{sec:vol9_canonical_source_routing}

This chapter is a synthesis chapter: no substrate primitive is re-derived. The canonical derivations live at:

\begin{itemize}
    \item \textbf{Axiom 4 canonical statement:} \kbleaf{manuscript/common\_equations/eq\_axiom\_4.tex} (single source of truth; INVARIANT-S2 mirror at \kbleaf{manuscript/ave-kb/CLAUDE.md}).
    \item \textbf{Universal saturation-kernel catalog:} \kbleaf{manuscript/ave-kb/common/universal-saturation-kernel-catalog.md} (\texttt{clm-gz7ryg}, \texttt{clm-dxdsvt}, \texttt{clm-hvvvop}, \texttt{clm-5fu303}, \texttt{clm-l4o7hv}; cross-scale catalog of record).
    \item \textbf{Four-regime map:} \kbleaf{manuscript/ave-kb/vol1/operators-and-regimes/ch7-regime-map/four-regimes.md} (\texttt{clm-2dwzib}, \texttt{clm-b2anl4}; regime-boundary derivations from first principles).
    \item \textbf{$\alpha$-invariance under SYM-class:} \kbleaf{manuscript/ave-kb/vol3/gravity/ch01-gravity-yield/alpha-invariance-symmetric-gravity.md} (\texttt{clm-3zz0f6}; canonical proof).
    \item \textbf{Two effective wave speeds ($c_{EM}$ vs $c_{shear}$):} \texttt{clm-8nkvwy} lines 111--113 at the same gravity-yield leaf above.
    \item \textbf{PONDER-05 substrate-physics anchor:} \kbleaf{manuscript/ave-kb/vol4/circuit-theory/ch1-vacuum-circuit-analysis/parametric-coupling-kernel.md} (Phase 0c + Phase 2-NA $P(\delta V)$ shape function derivation).
    \item \textbf{PONDER-05 operating-point canonical statement:} \kbleaf{manuscript/ave-kb/vol4/falsification/ch11-experimental-bench/measurement-hierarchy-snr.md}:66.
    \item \textbf{EE saturation translation:} \kbleaf{manuscript/ave-kb/common/translation-tables/translation-circuit.md} \S 9 (\texttt{clm-eemap1} META framework).
    \item \textbf{Vol~4 dielectric thresholds:} Vol~4 canonical leaf \texttt{clm-0vxzfu} ($V_{yield}$, $V_{snap}$).
    \item \textbf{Substrate-distinct $\alpha$-modulation under ASYM-class:} \kbleaf{ave-kb/vol3/cosmology/ch05-dark-sector/delta-strain-cosmic-tcc.md} (\texttt{clm-hp7nlm}, Cosserat-rotation-sector mass-gap thermal-mode-population ASYM mechanism).
\end{itemize}

\noindent
The natural vacuum substrate carries one saturation mechanism; this chapter documents its measurable characteristic curves; AVE substrate-physics derives the kernel from Axiom~4 alone; engineering practice empirically measures its dielectric, magnetic, semiconductor, and avalanche specializations at the corresponding component scales.
